%% ============================================================
%%  Lunar Thermal Model -- Plain-Language Technical Summary
%%  Paste into Overleaf and compile with pdfLaTeX
%% ============================================================
\documentclass[12pt, a4paper]{article}

\usepackage[margin=2.5cm]{geometry}
\usepackage{amsmath, amssymb}
\usepackage{booktabs}
\usepackage{xcolor}
\usepackage{hyperref}
\usepackage{parskip}
\usepackage{enumitem}
\usepackage{fancyhdr}
\usepackage{titlesec}
\usepackage{mdframed}
\usepackage{array}
\usepackage{microtype}

%% Colours
\definecolor{moonblue}{RGB}{25, 80, 150}
\definecolor{moongray}{RGB}{75, 75, 85}
\definecolor{moonwarm}{RGB}{170, 70, 20}
\definecolor{boxbg}{RGB}{238, 244, 255}
\definecolor{eqbg}{RGB}{255, 252, 230}

%% Hyperlinks
\hypersetup{colorlinks=true, linkcolor=moonblue,
            urlcolor=moonblue, citecolor=moonblue}

%% Section headings
\titleformat{\section}{\large\bfseries\color{moonblue}}{}{0em}{}[\titlerule]
\titleformat{\subsection}{\normalsize\bfseries\color{moongray}}{}{0em}{}
\titleformat{\subsubsection}{\normalsize\itshape\color{moongray}}{}{0em}{}

%% Header / footer
\pagestyle{fancy}\fancyhf{}
\rhead{\small Lunar Thermal Model Summary}
\lhead{\small \textcolor{moongray}{Gregorio}}
\rfoot{\small \thepage}

%% Blue info box
\newmdenv[backgroundcolor=boxbg, linecolor=moonblue, linewidth=1pt,
          roundcorner=4pt, innerleftmargin=10pt, innerrightmargin=10pt,
          innertopmargin=7pt, innerbottommargin=7pt]{plainbox}

%% Yellow equation-explanation box
\newmdenv[backgroundcolor=eqbg, linecolor=moonwarm, linewidth=1pt,
          roundcorner=4pt, innerleftmargin=10pt, innerrightmargin=10pt,
          innertopmargin=7pt, innerbottommargin=7pt]{eqbox}

%% ============================================================
\begin{document}

%% ---------- Title page --------------------------------------
\begin{titlepage}
  \centering\vspace*{1.8cm}
  {\Huge\bfseries\color{moonblue}
    One-Dimensional Lunar\\[6pt] Subsurface Thermal Model}\\[1.2cm]
  {\large A Plain-Language Technical Summary\\
   from First Principles to Python Results}\\[2.2cm]
  {\large\itshape Gregorio, R.\,P.}\\[0.5cm]
  {\normalsize\color{moongray}
   Based on the \texttt{Lunar-Clean} Python notebook\\[2pt]
   \url{https://github.com/rp3gregorio/Lunar-Clean}}\\[3cm]
  {\large \today}
  \vfill
  \begin{plainbox}
    \textbf{Purpose of this document.}
    This paper explains, in plain English, every equation and every
    number used in the lunar thermal model.
    No prior background in physics or planetary science is assumed.
    After reading it you should be able to say:
    \emph{``I know what the model is doing, why each number is there,
    and what the outputs mean.''}
  \end{plainbox}
\end{titlepage}

\tableofcontents
\newpage

%% ============================================================
\section{Introduction -- Why Does the Moon Get So Hot and Cold?}

The Moon has no atmosphere.
On Earth the air acts like a blanket: it slows down how fast heat
arrives from the Sun and how fast it escapes at night.
On the Moon there is no blanket.
As a result the surface swings from roughly
\textbf{100\,K} ($-173\,^{\circ}$C) at night to
\textbf{390\,K} ($+117\,^{\circ}$C) in the middle of the day --
a swing of nearly 300 degrees, about five times the most extreme
range ever recorded on Earth.

A \textbf{lunar day} lasts 29.53 Earth days (one full rotation of
the Moon relative to the Sun).
The surface bakes for $\sim$two Earth weeks, then freezes for
$\sim$two more weeks.

This model answers: \emph{how does that heat wave travel downward
into the soil?}
By the time you are 1\,metre underground the temperature barely
changes at all -- the soil insulates itself.
Understanding this depth profile matters for:

\begin{itemize}[leftmargin=1.5em, itemsep=2pt]
  \item Designing future lunar mission hardware (batteries,
        electronics, and human habitats all have temperature limits).
  \item Understanding the Moon's internal heat and geological history.
  \item Interpreting satellite observations from the Diviner radiometer
        aboard the Lunar Reconnaissance Orbiter (LRO).
\end{itemize}

%% ============================================================
\section{The Lunar Soil (Regolith)}

The top few metres of the Moon are covered by a layer of fine,
powdery, pulverised rock called \textbf{regolith}.
It has been ground up by billions of years of meteorite impacts.
Its thermal properties are very different from ordinary rock:

\begin{itemize}[leftmargin=1.5em, itemsep=2pt]
  \item It is \textbf{extremely fluffy near the surface} --
        imagine very dry, loose talcum powder.
        This makes it a poor conductor of heat.
  \item It becomes \textbf{progressively denser and more compact
        with depth} because the weight of overlying material
        compresses it.
  \item Apollo astronauts drilled into it to about 2\,metres and
        measured temperature at several depths.
        These measurements are used to validate the model.
\end{itemize}

%% ============================================================
\section{The Digital Elevation Model (DEM)}

Before simulating temperature the model needs to know the shape of
the terrain.
It uses the \textbf{LOLA LDEM} -- the Lunar Orbiter Laser Altimeter
Global Digital Elevation Model -- produced by NASA's LRO mission.

\begin{plainbox}
  \textbf{What a DEM is.}
  Think of it as a huge grid of altitude measurements covering
  the whole Moon.
  Each grid cell holds one number: the height of that spot above
  a reference sphere of radius
  $R_{\Moon} = 1{,}737{,}400$\,m (the Moon's mean radius).
\end{plainbox}

The version included in this repository (\texttt{LDEM\_4}) has:

\begin{center}
\begin{tabular}{ll}
\toprule
Property & Value \\
\midrule
Grid size       & $720 \times 1440$ pixels \\
Resolution      & 4 pixels per degree \\
Pixel size      & $\approx 7{,}581$\,m per pixel \\
Elevation range & $-17{,}758$\,m to $+21{,}008$\,m \\
Elevation unit  & metres relative to $R_{\Moon}$ \\
Digital number scale & 1\,DN $= 0.5$\,m \\
\bottomrule
\end{tabular}
\end{center}

The DEM is used for three purposes:
(1) reading the elevation, slope, and aspect of the target location;
(2) computing the \textbf{horizon profile} (how high surrounding
hills appear); and
(3) drawing the terrain maps in Section~2b of the notebook.

\subsection{Converting coordinates to pixels}

The DEM uses a \textbf{simple cylindrical} (equirectangular)
projection, with column~0 at $0^{\circ}$E and row~0 at $90^{\circ}$N.
To convert latitude $\phi$ and longitude $\lambda$ to a pixel:

\begin{equation}
  \text{row} =
    \left\lfloor \frac{90 - \phi}{\,p\,} - 0.5 \right\rceil,
  \qquad
  \text{col} =
    \left\lfloor \frac{\lambda}{\,p\,} - 0.5 \right\rceil
  \label{eq:pixel}
\end{equation}

where $p = 1/\text{map\_res} = 0.25^{\circ}$/pixel is the angular
pixel size and $\lfloor\cdot\rceil$ means ``round to nearest integer''.

\subsection{Slope and aspect from the DEM}

The \textbf{slope} $\alpha$ is the steepness of the terrain
(degrees from horizontal).
The \textbf{aspect} $\beta$ is the compass direction the slope
faces ($0 = $ north, $90^{\circ} = $ east).
Both are computed by taking height differences between neighbouring
pixels (\emph{central differences}):

\begin{equation}
  \frac{\partial z}{\partial x}
    \approx \frac{z_{i,\,j+1}-z_{i,\,j-1}}{2\,\Delta x},
  \qquad
  \frac{\partial z}{\partial y}
    \approx \frac{z_{i+1,\,j}-z_{i-1,\,j}}{2\,\Delta y}
  \label{eq:gradient}
\end{equation}

\begin{equation}
  \alpha = \arctan\!\left(
    \sqrt{\left(\frac{\partial z}{\partial x}\right)^2 +
          \left(\frac{\partial z}{\partial y}\right)^2}
  \right),
  \qquad
  \beta  = \arctan\!\left(
    \frac{\partial z/\partial x}{-\partial z/\partial y}
  \right)
  \label{eq:slope_aspect}
\end{equation}

where $\Delta x = \Delta y = \texttt{pixel\_m} \approx 7{,}581$\,m.

%% ============================================================
\section{Solar Geometry -- Where Is the Sun?}

The amount of energy hitting the surface at any moment depends
entirely on the Sun's position in the sky.

\subsection{Lunar rotation and the solar day}

The Moon completes one rotation relative to the Sun in one synodic
month:

\begin{equation}
  T_{\text{lunar}} = 29.53 \times 86{,}400
                   \approx 2{,}551{,}392\,\text{s}
  \label{eq:lunar_day}
\end{equation}

\subsection{Solar declination (seasonal wobble)}

The Moon's spin axis is tilted by only $1.54^{\circ}$ relative to
its orbital plane (Earth's tilt is $23.5^{\circ}$).
This tiny tilt produces a small seasonal variation in the Sun's
angle above the equator:

\begin{equation}
  \delta(t) = 1.54^{\circ} \times
  \sin\!\left(\frac{2\pi\,t}{27.3 \times 86{,}400}\right)
  \label{eq:declination}
\end{equation}

where $t$ is time in seconds since an arbitrary epoch.

\begin{plainbox}
  \textbf{Plain English.}
  $\delta(t)$ is the angle between the Sun and the lunar equator.
  Because it is so small ($\pm 1.54^{\circ}$), the Moon has almost
  no seasons.
  The dominant effect is the daily (monthly) rotation -- local noon
  versus local midnight.
\end{plainbox}

\subsection{Hour angle}

The \textbf{hour angle} $h$ describes the Sun's east--west position
in the sky; it increases as the Moon rotates:

\begin{equation}
  h(t) = \frac{2\pi\,t}{T_{\text{lunar}}} + \lambda_{\text{rad}}
  \label{eq:hour_angle}
\end{equation}

where $\lambda_{\text{rad}}$ is the site longitude in radians.
When $h = 0$ the Sun is directly overhead for an equatorial site.

\subsection{Solar zenith angle}

The \textbf{zenith angle} $\theta_z$ measures how far the Sun is
from being directly overhead ($\theta_z = 0^{\circ}$ = overhead,
$\theta_z = 90^{\circ}$ = on the horizon):

\begin{equation}
  \cos\theta_z
  = \sin\phi\,\sin\delta + \cos\phi\,\cos\delta\,\cos h
  \label{eq:zenith}
\end{equation}

where $\phi$ is the site latitude.

\begin{plainbox}
  \textbf{Intuition.}
  On the equator at local noon ($h=0$, $\delta\approx0$):
  $\cos\theta_z = 1$, so the Sun is directly overhead.
  At the poles the Sun skims the horizon at most times, which is
  why some polar craters remain permanently in shadow.
\end{plainbox}

\subsection{Absorbed solar flux on a sloped surface}

The total solar energy available at the Moon is the
\textbf{solar constant}:
\begin{equation}
  S_0 = 1361\,\text{W\,m}^{-2}
  \label{eq:solar_constant}
\end{equation}

Not all of it is absorbed.
The \textbf{albedo} $A$ is the fraction reflected away
(lunar regolith: $A \approx 0.09$, so 91\% is absorbed).
A sloped surface tilted at angle $\alpha$ and facing direction
$\beta$ intercepts a different amount of sunlight than a flat surface.
The absorbed solar flux is:

\begin{equation}
  Q_{\odot} = S_{\text{scale}}\,(1 - A)\,S_0\,\cos\theta_i
  \label{eq:solar_flux}
\end{equation}

where the \textbf{incidence angle} $\theta_i$ between the Sun
direction and the surface normal is:

\begin{equation}
  \cos\theta_i =
    \cos\theta_z\cos\alpha
    + \sin\theta_z\sin\alpha\cos(\beta - \theta_{az})
  \label{eq:incidence}
\end{equation}

and $\theta_{az}$ is the solar azimuth.
The factor $S_{\text{scale}}$ (``SUNSCALE'', default 1.10) is
an adjustable calibration multiplier.
If $\cos\theta_i \le 0$ the surface faces away from the Sun
and $Q_{\odot} = 0$.

%% ============================================================
\section{Topographic Shadowing}

A location can receive zero sunlight even when the Sun is
geometrically above the horizon if a nearby hill blocks it.

\subsection{Horizon elevation profile}

For each of $N_{\varphi} = 360$ azimuth directions $\varphi_k$,
the model scans outward from the target pixel and records the
maximum angular elevation of any terrain obstacle:

\begin{equation}
  H(\varphi_k)
  = \max_{d}\; \arctan\!\left(
      \frac{z(d,\varphi_k) - z_0}{\,d\cdot\Delta x\,}
    \right)
  \label{eq:horizon}
\end{equation}

where $z_0$ is the target elevation and $d$ is the horizontal
distance to the obstacle pixel.

\subsection{Illumination check}

The surface is lit if and only if the Sun's elevation angle
exceeds the terrain horizon in the Sun's direction:

\begin{equation}
  \text{lit} =
  \begin{cases}
    \mathrm{True}  &
      \text{if } \theta_z < \tfrac{\pi}{2}
      \text{ and }
      \bigl(\tfrac{\pi}{2} - \theta_z\bigr) > H(\theta_{az}) \\[4pt]
    \mathrm{False} & \text{otherwise}
  \end{cases}
  \label{eq:illumination}
\end{equation}

\begin{plainbox}
  \textbf{Plain English.}
  Imagine shining a torch toward a hill.
  The torch (Sun) must be higher in the sky than the top of the
  hill for the light to reach the ground in front of you.
  The horizon profile records the ``top of every hill'' in every
  compass direction.
\end{plainbox}

\subsection{Sky-view factor (SVF)}

The \textbf{sky-view factor} is a single number summarising how
much of the open sky hemisphere is visible from the target point.
SVF = 1.0 means flat, unobstructed terrain; SVF $< 1$ means
surrounding terrain blocks part of the sky.

\begin{equation}
  \mathrm{SVF}
  = \frac{1}{2\pi}\int_0^{2\pi}\cos^3\!\bigl(H(\varphi)\bigr)\,
    \mathrm{d}\varphi
  \approx
  \frac{1}{N_{\varphi}}\sum_{k=1}^{N_{\varphi}}
    \cos^3\!\bigl(H(\varphi_k)\bigr)
  \label{eq:svf}
\end{equation}

%% ============================================================
\section{The Heat Conduction Equation -- The Core of the Model}

\subsection{Physical idea}

When one part of a material is hotter than a neighbouring part,
heat flows from hot to cold.
The rate of flow depends on how easily the material conducts heat
(its \emph{thermal conductivity}).
In the Moon's regolith this process describes how the surface
warmth (or cold) spreads downward into the soil.

\subsection{Governing equation}

The model solves the \textbf{one-dimensional heat conduction
equation}:

\begin{equation}
  \rho(z)\,c(T)\,\frac{\partial T}{\partial t}
  =
  \frac{\partial}{\partial z}
  \!\left[k(T,z)\,\frac{\partial T}{\partial z}\right]
  \label{eq:heat_eq}
\end{equation}

\begin{eqbox}
  \textbf{Reading the equation term by term.}
  \begin{itemize}[leftmargin=1.2em, itemsep=3pt]
    \item $T(z,t)$ -- temperature (K) at depth $z$ (m) and
          time $t$ (s).
    \item $\rho(z)$ -- density (kg/m$^3$) of the soil at depth $z$.
    \item $c(T)$ -- specific heat capacity (J/(kg\,K)):
          energy needed to raise 1\,kg by 1\,K.
    \item \textbf{Left side} $\rho c\,\partial T/\partial t$:
          rate at which energy is stored in the soil.
    \item $k(T,z)$ -- thermal conductivity (W/(m\,K)):
          how easily heat flows through the material.
    \item $k\,\partial T/\partial z$ -- heat flux (W/m$^2$)
          flowing vertically.
    \item \textbf{Right side}: net heat arriving at each point
          from its neighbours above and below.
  \end{itemize}
\end{eqbox}

%% ============================================================
\section{Material Properties of Lunar Regolith}

\subsection{Thermal conductivity}

Conductivity in regolith has two components:

\begin{enumerate}[leftmargin=1.5em, itemsep=2pt]
  \item \textbf{Solid (contact) conduction} $k_s(z)$ -- heat
        transferred through direct grain-to-grain contact.
        Depends on depth because density increases with depth.
  \item \textbf{Radiative conduction} -- at high temperatures
        grains exchange infrared radiation with each other.
        This component grows steeply with temperature.
\end{enumerate}

The total conductivity is:

\begin{equation}
  k(T,z) = k_s(z)\,\left[1 + \chi\left(\frac{T}{T_{\mathrm{ref}}}\right)^3\right]
  \label{eq:conductivity}
\end{equation}

where $T_{\mathrm{ref}} = 350$\,K and $\chi = 2.7$ is the
\textbf{radiative conductivity parameter}
(adjustable; higher $\chi$ means more radiative heat transfer).
The $T^3$ dependence comes directly from Stefan-Boltzmann
blackbody physics.

\begin{plainbox}
  \textbf{Why conductivity matters so much.}
  The contact conductivity at the surface is
  $k_s \approx 7\times10^{-4}$\,W/(m\,K) -- roughly 1000 times
  less than solid rock.
  This extreme insulation is why the surface temperature swings
  so violently: heat cannot escape downward fast enough to damp
  out the day--night cycle.
\end{plainbox}

\subsection{Specific heat capacity}

Heat capacity $c(T)$ is how much energy (joules) it takes to heat
1\,kg of regolith by 1\,K.
It rises with temperature.
The model uses the polynomial fit of Hemingway et al.\ (1973),
measured on actual Apollo lunar samples:

\begin{equation}
  c(T) = -3.6125
       + 2.7431\,T
       + 2.3616\times10^{-3}\,T^2
       - 1.234\times10^{-5}\,T^3
       + 8.9093\times10^{-9}\,T^4
  \quad\bigl[\mathrm{J\,kg^{-1}\,K^{-1}}\bigr]
  \label{eq:heat_capacity}
\end{equation}

Typical values:
$c(100\,\mathrm{K}) \approx 270$\,J/(kg\,K) at night;
$c(390\,\mathrm{K}) \approx 1050$\,J/(kg\,K) at peak day.

\subsection{Density models}

\subsubsection{Model 1 -- Discrete layers (Apollo drill-core)}

Apollo astronauts observed three distinct compaction layers.
The boundary between Layer~1 and Layer~2 is $H$ (default 7\,cm):

\begin{equation}
  \rho_{\mathrm{disc}}(z) =
  \begin{cases}
    \rho_s
      & 0 \le z < H \\[4pt]
    \rho_s + \dfrac{1700 - \rho_s}{L_2 - H}\,(z - H)
      & H \le z < L_2 \\[8pt]
    1700 + 100\,\min\!\left(1,\;\dfrac{z - L_2}{2.80}\right)
      & z \ge L_2
  \end{cases}
  \label{eq:density_discrete}
\end{equation}

where $\rho_s = 1100$\,kg/m$^3$ (surface density, adjustable),
$L_2 = 0.20$\,m (20\,cm, fixed Layer~2/3 boundary).
The Layer~1 depth $H \approx 7$\,cm and the Layer~2 reference density
1700\,kg/m$^3$ are based on Apollo drill-core measurements
(Carrier et al.\ 1991).

\begin{center}
\begin{tabular}{llll}
\toprule
Layer & Depth & Density & Description \\
\midrule
1 & 0 to $H$ ($\approx$7\,cm) & 1100\,kg/m$^3$ & Fluffy surface dust \\
2 & $H$ to 20\,cm & 1100--1700\,kg/m$^3$ & Transitional (linear ramp) \\
3 & $>$20\,cm & 1700--1800\,kg/m$^3$ & Consolidated regolith \\
\bottomrule
\end{tabular}
\end{center}

\subsubsection{Model 2 -- Hayne et al.\ 2017 (exponential compaction)}

Diviner satellite data suggest a smooth exponential increase
in density with depth (Hayne et al.\ 2017):

\begin{equation}
  \rho_H(z) = \rho_d - (\rho_d - \rho_s)\,
              \exp\!\left(-\frac{z}{H}\right)
  \label{eq:density_hayne}
\end{equation}

with $\rho_s = 1100$\,kg/m$^3$, $\rho_d = 1800$\,kg/m$^3$,
and $H = 0.07$\,m.
Note that the global best-fit from Hayne et al.\ (2017) is $H = 0.06$\,m;
we adopt $H = z_1 = 0.07$\,m so that both models share the same
characteristic compaction depth, isolating the effect of density
parameterisation rather than the choice of depth scale.

Solid conductivity follows the same exponential shape:

\begin{equation}
  k_s^H(z) = k_{\mathrm{surf}}
    + (k_{\mathrm{deep}} - k_{\mathrm{surf}})\,
      \frac{\rho_H(z) - \rho_s}{\rho_d - \rho_s}
  \label{eq:k_hayne}
\end{equation}

with $k_{\mathrm{surf}} = 7.4\times10^{-4}$\,W/(m\,K) and
$k_{\mathrm{deep}} = 3.8\times10^{-3}$\,W/(m\,K) \citep{Hayne2017}.

%% ============================================================
\section{Boundary Conditions}

The heat equation needs to know what happens at the top and
bottom of the soil column.

\subsection{Top boundary -- surface energy balance}

At the surface ($z = 0$), energy arrives from the Sun and leaves
as infrared radiation.
These must balance:

\begin{equation}
  \underbrace{k(T_0,0)\,\frac{T_1 - T_0}{\Delta z}}_{\text{conduction up}}
  +\;
  \underbrace{Q_{\odot}}_{\text{solar in}}
  =\;
  \underbrace{\varepsilon\,\sigma\,T_0^4}_{\text{IR radiation out}}
  \label{eq:surface_bc}
\end{equation}

where:
\begin{itemize}[leftmargin=1.5em, itemsep=2pt]
  \item $T_0$ -- unknown surface temperature (K)
  \item $T_1$ -- temperature at the first sub-surface grid node
  \item $\Delta z$ -- spacing between nodes (5\,mm)
  \item $\varepsilon = 0.95$ -- \textbf{emissivity} (how efficiently
        the surface radiates as a blackbody)
  \item $\sigma = 5.6704\times10^{-8}$\,W\,m$^{-2}$\,K$^{-4}$ --
        Stefan-Boltzmann constant
\end{itemize}

\begin{plainbox}
  \textbf{Plain English.}
  The surface temperature adjusts itself until the heat it radiates
  away exactly balances the heat arriving from the Sun plus any
  heat conducted up from below.
  On the dayside the Sun dominates; at night the surface cools
  rapidly because it keeps radiating but receives nothing back.
\end{plainbox}

Because of the $T_0^4$ term, equation~\eqref{eq:surface_bc} is
nonlinear and cannot be solved directly.
The model uses the \textbf{Newton-Raphson method}: an iterative
approach that converges in typically 3--5 steps.

\subsection{Bottom boundary -- basal heat flux}

At 3\,m depth a small amount of heat flows upward from the Moon's
hot interior.
This was measured directly by the Apollo Heat Flow Experiment:

\begin{equation}
  Q_{\mathrm{basal}} = 18\,\mathrm{mW\,m}^{-2}
  \label{eq:basal}
\end{equation}

Implemented as a Neumann (flux) condition:

\begin{equation}
  T_{\mathrm{bottom}} = T_{\mathrm{bottom}-1}
    + \frac{Q_{\mathrm{basal}}\,\Delta z_{\mathrm{bottom}}}{k_{\mathrm{bottom}}}
  \label{eq:bottom_bc}
\end{equation}

\begin{plainbox}
  \textbf{Why 18\,mW/m$^2$ matters.}
  This is tiny compared to peak solar flux ($\sim$1200\,W/m$^2$),
  so it has almost no effect on surface temperatures.
  However, it sets the temperature at depth ($>$1\,m), where
  Apollo drills measured $\sim$253--257\,K.
  Getting this right confirms that the deep physics is correct.
\end{plainbox}

%% ============================================================
\section{Numerical Method -- Solving the Equations on a Computer}

\subsection{The depth grid}

The soil column is divided into discrete \textbf{nodes}
(sample points at fixed depths):

\begin{center}
\begin{tabular}{lll}
\toprule
Region & Spacing & Reason \\
\midrule
0 to 10\,cm & 5\,mm & Steep temperature gradients near surface \\
10\,cm to 3\,m & 20\,mm & Small gradients; coarser grid saves time \\
\bottomrule
\end{tabular}
\end{center}

This gives approximately 70--80 nodes total.

\subsection{Explicit finite-difference time-stepping}

At each interior node $i$, temperature is updated one time step
at a time using the \textbf{explicit forward-Euler} scheme:

\begin{equation}
  T_i^{n+1} = T_i^n +
  \frac{\Delta t}{\rho_i\,c_i\,\overline{\Delta z}_i}
  \left[
    k_{i+\frac{1}{2}}\,\frac{T_{i+1}^n - T_i^n}{\Delta z_i}
    -
    k_{i-\frac{1}{2}}\,\frac{T_i^n - T_{i-1}^n}{\Delta z_{i-1}}
  \right]
  \label{eq:fd}
\end{equation}

Superscripts $n$ and $n+1$ denote the current and next time step.
Interface conductivities are arithmetic means:
$k_{i+1/2} = \tfrac{1}{2}(k_i + k_{i+1})$.

\subsection{Stability criterion}

The explicit scheme is only stable when the time step is small
enough.
The \textbf{CFL stability condition} requires:

\begin{equation}
  \Delta t < \frac{1}{2}
  \frac{\rho_{\max}\,c_{\max}\,(\Delta z_{\min})^2}{k_{\max}}
  \label{eq:stability}
\end{equation}

The model uses a safety factor $f = 0.01$, giving
$\Delta t \approx 100$\,s.
Running 3 lunar days therefore requires $\sim$76{,}000 time steps
(and takes roughly 15--30 seconds on a modern laptop).

\subsection{Spin-up period}

The simulation starts from a uniform temperature of 250\,K.
After $\sim$2 lunar days the temperature field reaches a
\textbf{periodic steady state} -- it repeats the same pattern
every lunar day.
The model runs for $N_{\mathrm{days}}$ days (default 3) and
analysis is performed only on the \textbf{final} day,
discarding the earlier spin-up.

%% ============================================================
\section{Complete Parameter Reference}

\begin{center}
\renewcommand{\arraystretch}{1.25}
\begin{tabular}{>{\ttfamily}l l l p{5.5cm}}
\toprule
\textrm{Name} & Symbol & Value & Meaning \\
\midrule
\multicolumn{4}{l}{\itshape Physical constants (never change)} \\[2pt]
S0         & $S_0$     & 1361\,W\,m$^{-2}$               & Solar constant at 1\,AU \\
sigma      & $\sigma$  & $5.6704\!\times\!10^{-8}$\,W\,m$^{-2}$\,K$^{-4}$ & Stefan-Boltzmann \\
Q\_basal   & $Q_b$     & 0.018\,W\,m$^{-2}$              & Basal heat flux (Apollo HFE) \\
R\_MOON    & $R_{\Moon}$& 1{,}737{,}400\,m               & Moon mean radius \\
LUNAR\_DAY & $T_{\odot}$& 2{,}551{,}392\,s               & 29.53 Earth days \\[6pt]
\multicolumn{4}{l}{\itshape Adjustable model parameters} \\[2pt]
SUNSCALE   & $S_{\rm sc}$ & 1.10                        & Solar flux multiplier \\
ALBEDO     & $A$          & 0.09                        & Fraction of sunlight reflected \\
EMISSIVITY & $\varepsilon$& 0.95                        & IR radiation efficiency \\
CHI        & $\chi$       & 2.7                         & Radiative conductivity exponent \\
H\_PARAM   & $H$          & 0.07\,m                     & Layer-1 thickness / scale height (Apollo drill data; Carrier et al.\ 1991) \\
NDAYS      & --           & 3                           & Lunar days to simulate \\
T\_INIT    & $T_0$        & 250\,K                      & Initial uniform temperature \\[6pt]
\multicolumn{4}{l}{\itshape Regolith properties} \\[2pt]
rho\_s     & $\rho_s$     & 1100\,kg\,m$^{-3}$          & Surface density (adjustable) \\
rho\_d     & $\rho_d$     & 1800\,kg\,m$^{-3}$          & Deep density (Carrier et al.\ 1991) \\
k\_surf    & $k_{\rm surf}$& $7.4\!\times\!10^{-4}$\,W/(m\,K) & Surface conductivity \\
k\_deep    & $k_{\rm deep}$& $3.8\!\times\!10^{-3}$\,W/(m\,K) & Deep conductivity \citep{Hayne2017} \\
T\_ref     & $T_{\rm ref}$ & 350\,K                    & Reference T for $\chi$ term \\[6pt]
\multicolumn{4}{l}{\itshape Numerical grid} \\[2pt]
DZ\_FINE   & $\Delta z_f$  & 5\,mm    & Node spacing near surface \\
DZ\_COARSE & $\Delta z_c$  & 20\,mm   & Node spacing at depth \\
DEPTH\_FINE& --            & 0.10\,m  & Switch depth fine/coarse \\
DEPTH\_MAX & --            & 3.0\,m   & Total model depth \\
\bottomrule
\end{tabular}
\end{center}

%% ============================================================
\section{Apollo Site Validation}

The most important test of any model is whether it reproduces
real measurements.
Two Apollo missions placed thermometers at multiple depths in
the regolith:

\begin{center}
\begin{tabular}{lllr}
\toprule
Mission & Coordinates & Drill depth & Measured T at depth \\
\midrule
Apollo 15 (1971) & 26.13°N, 3.63°E  & $\sim$2.0\,m & 252.8--253.3\,K \\
Apollo 17 (1972) & 20.19°N, 30.77°E & $\sim$2.3\,m & 255.5--256.9\,K \\
\bottomrule
\end{tabular}
\end{center}

\begin{plainbox}
  \textbf{Why is Apollo 17 slightly warmer underground?}
  It sits closer to the equator ($20.2^{\circ}$ vs $26.1^{\circ}$N)
  and therefore receives more sunlight on average throughout the
  year.
  More solar input warms the deep soil.
  At $\sim$2\,m depth the temperature no longer swings with the
  day--night cycle; it reflects only the \emph{long-term average}
  solar heating, which is latitude-dependent.
\end{plainbox}

The model accuracy is reported using two metrics:

\begin{align}
  \mathrm{RMSE} &= \sqrt{\frac{1}{N}\sum_{i=1}^{N}
    \bigl(T_{\mathrm{model}}(z_i) - T_{\mathrm{Apollo}}(z_i)\bigr)^2}
  \label{eq:rmse}\\[6pt]
  \mathrm{Bias} &= \frac{1}{N}\sum_{i=1}^{N}
    \bigl(T_{\mathrm{model}}(z_i) - T_{\mathrm{Apollo}}(z_i)\bigr)
  \label{eq:bias}
\end{align}

A well-calibrated model achieves RMSE $< 1$\,K.
Positive bias means the model runs hot; negative bias means cold.

%% ============================================================
\section{Guide to the Notebook Sections}

\begin{center}
\renewcommand{\arraystretch}{1.3}
\begin{tabular}{>{\bfseries}c l p{8.2cm}}
\toprule
\S & Title & What to look for \\
\midrule
1  & Configuration         & All adjustable parameters live here. Edit this cell, then re-run Sections 2 onward. \\
1b & Parameter Panel       & Sliders -- change parameters without editing any code; click Apply then re-run Section 3. \\
2  & Load DEM              & Confirms the DEM file was found. Prints elevation, slope, sky-view factor of the target pixel. \\
2b & Terrain Map           & Red star = target location. Contour lines show terrain shape around it. \\
2c & Horizon Profile       & Tall brown wedges in a direction = a big hill blocking the sky there. A flat ring = open terrain. \\
2d & Illumination Timeline & Gold = sunlit; grey = in topographic shadow; blue = night. Width of the gold band = length of the lunar daytime at this site. \\
2e & Temperature Map       & Centre panel: estimated peak noon temperature. Right panel: day-night swing -- largest near the equator. \\
3  & Run Model             & Runs the full simulation ($\sim$15--30\,s). Prints T range. If T\textsubscript{max} $>$ 400\,K the model runs hot -- try lowering SUNSCALE. \\
4  & Diurnal Cycles        & Shows how quickly the surface heats and cools and how the heat wave damps out with depth. \\
5  & Heatmap               & Red = hot (day), black = cold (night). Notice the pattern fades below $\sim$20\,cm depth. \\
6  & Apollo Comparison     & Green dots = measured; coloured line = model. RMSE $< 1$\,K is an excellent fit. \\
7  & Model Comparison      & How much does the choice of density model (discrete vs Hayne) matter? \\
8  & Sensitivity Analysis  & Which parameter has the biggest impact on accuracy? Options include \texttt{rho\_surface}. \\
9  & Single Point          & Run two different models at any list of locations side by side. \\
10 & Batch Processing      & Process many locations at once and compare their surface temperatures. \\
11 & Density Model Viewer  & Move the slider; the density-vs-depth profile redraws instantly. \\
\bottomrule
\end{tabular}
\end{center}

%% ============================================================
\section{Tuning Guide -- When the Model is Off}

\begin{center}
\renewcommand{\arraystretch}{1.3}
\begin{tabular}{p{4.5cm} p{4.5cm} p{4.5cm}}
\toprule
\textbf{Symptom} & \textbf{Likely cause} & \textbf{Fix} \\
\midrule
Surface T\textsubscript{max} too low ($<$360\,K at equator) &
  Not enough solar energy &
  Increase SUNSCALE (try 1.05--1.15) or decrease ALBEDO \\
Surface T\textsubscript{max} too high &
  Too much energy absorbed &
  Decrease SUNSCALE or increase ALBEDO \\
Subsurface T too low (negative bias vs Apollo) &
  Not enough mean energy input &
  Increase SUNSCALE slightly \\
Diurnal swing too large at depth &
  Conductivity too low &
  Increase CHI \\
Density profile unrealistic &
  H\_PARAM or rho\_surface wrong &
  Adjust via Section 11 sliders, then re-run \\
\bottomrule
\end{tabular}
\end{center}

%% ============================================================
\section{End-to-End Model Flow}

The complete simulation runs in this order:

\begin{enumerate}[leftmargin=1.8em, itemsep=5pt]
  \item \textbf{Load DEM.}
        Read the elevation grid (\texttt{LDEM\_4.IMG}).
        Convert the target (lat, lon) to a pixel using
        equation~\eqref{eq:pixel}.
        Compute slope $\alpha$ and aspect $\beta$
        (equations \ref{eq:gradient}--\ref{eq:slope_aspect}).

  \item \textbf{Compute horizon profile.}
        Scan outward in $N_{\varphi}=360$ directions.
        Record $H(\varphi_k)$ at each azimuth
        (equation~\ref{eq:horizon}).
        Compute SVF (equation~\ref{eq:svf}).

  \item \textbf{Initialise.}
        Set $T(z,0) = 250$\,K everywhere.
        Build the non-uniform depth grid ($\sim$75 nodes,
        0 to 3\,m).

  \item \textbf{Time loop} (repeated $\sim$76{,}000 times for 3 days).
    \begin{enumerate}[leftmargin=1.4em, itemsep=3pt]
      \item Compute $\delta(t)$, $h(t)$, and $\cos\theta_z$
            (equations \ref{eq:declination}--\ref{eq:zenith}).
      \item Check illumination using $H(\varphi_k)$
            (equation~\ref{eq:illumination}).
      \item If lit: compute $Q_{\odot}$
            (equations \ref{eq:solar_flux}--\ref{eq:incidence}).
      \item Solve surface energy balance for $T_0$ via
            Newton-Raphson (equation~\ref{eq:surface_bc}).
      \item Update all interior nodes using finite differences
            (equation~\ref{eq:fd}).
      \item Apply bottom boundary condition
            (equation~\ref{eq:bottom_bc}).
    \end{enumerate}

  \item \textbf{Discard spin-up.}
        Keep only the final lunar day of data.

  \item \textbf{Post-process.}
        Extract T\textsubscript{min}, T\textsubscript{max},
        T\textsubscript{mean}, and amplitude at each depth.
        Compute RMSE and bias vs Apollo data if applicable.

  \item \textbf{Plot.}
        Generate all figures automatically.
\end{enumerate}

%% ============================================================
\section{References}

\begin{itemize}[leftmargin=1.5em, itemsep=5pt]

  \item Carrier, W.\,D., III, Olhoeft, G.\,R., \& Mendell, W.\ (1991).
        ``Physical Properties of the Lunar Surface.''
        In \textit{Lunar Sourcebook: A User's Guide to the Moon},
        G.\,H.\ Heiken, D.\,T.\ Vaniman, \& B.\,M.\ French (eds.),
        Cambridge University Press, pp.\ 475--594.

  \item Hayne, P.\,O., et al.\ (2017).
        ``Global Regolith Thermophysical Properties of the Moon
        from the Diviner Lunar Radiometer Experiment.''
        \textit{Journal of Geophysical Research: Planets},
        122, 2371--2400.
        \href{https://doi.org/10.1002/2017JE005387}
             {doi:10.1002/2017JE005387}

  \item Hemingway, B.\,S., Robie, R.\,A., \& Wilson, W.\,H.\
        (1973). ``Specific Heats of Lunar Soils, Basalt, and
        Breccias from the Apollo 14, 15, and 16 Landing Sites.''
        \textit{Proc.\ Fourth Lunar Science Conference}, 3,
        2481--2487.

  \item Langseth, M.\,G., Keihm, S.\,J., \& Peters, K.\ (1976).
        ``Revised Lunar Heat-Flow Values.''
        \textit{Proc.\ Seventh Lunar Science Conference},
        3143--3171.

  \item Smith, D.\,E., et al.\ (2017).
        LRO LOLA Global Digital Elevation Model V3.0.
        \textit{NASA Planetary Data System}.
        Dataset ID: \texttt{LRO-L-LOLA-4-GDR-V1.0}.

  \item Vasavada, A.\,R., et al.\ (2012).
        ``Lunar Equatorial Surface Temperatures and Regolith
        Properties from the Diviner Lunar Radiometer Experiment.''
        \textit{Journal of Geophysical Research: Planets},
        117, E00H18.

\end{itemize}

%% ============================================================
\appendix
\section{Quick-Reference: All Key Equations}

\renewcommand{\arraystretch}{1.6}
\begin{center}
\begin{tabular}{p{4.8cm} p{9cm}}
\toprule
\textbf{What it computes} & \textbf{Equation} \\
\midrule
Solar declination &
  $\delta = 1.54^{\circ}\sin\!\bigl(2\pi t/27.3\,\mathrm{days}\bigr)$ \\
Hour angle &
  $h = 2\pi t / T_{\odot} + \lambda_{\rm rad}$ \\
Solar zenith angle &
  $\cos\theta_z = \sin\phi\sin\delta + \cos\phi\cos\delta\cos h$ \\
Incidence on slope &
  $\cos\theta_i = \cos\theta_z\cos\alpha
   + \sin\theta_z\sin\alpha\cos(\beta-\theta_{az})$ \\
Absorbed solar flux &
  $Q_{\odot} = S_{\rm sc}(1-A)\,S_0\cos\theta_i$ \\
Hayne density &
  $\rho = 1800 - 700\exp(-z/H)$\quad [kg\,m$^{-3}$] \\
Total conductivity &
  $k = k_s(z)\bigl[1 + \chi(T/350)^3\bigr]$\quad [W\,m$^{-1}$K$^{-1}$] \\
Heat capacity &
  $c = -3.61 + 2.74T + \cdots$\quad [J\,kg$^{-1}$\,K$^{-1}$] \\
Heat conduction PDE &
  $\rho c\,\partial T/\partial t
   = \partial/\partial z\bigl[k\,\partial T/\partial z\bigr]$ \\
Surface energy balance &
  $k(T_1-T_0)/\Delta z + Q_{\odot} = \varepsilon\sigma T_0^4$ \\
Bottom BC &
  $T_N = T_{N-1} + Q_b\,\Delta z / k$ \\
Sky-view factor &
  $\mathrm{SVF} = N^{-1}\sum_k\cos^3 H(\varphi_k)$ \\
\bottomrule
\end{tabular}
\end{center}

\end{document}
